\documentclass[12pt,a4paper]{report}
\usepackage[utf8]{inputenc}
\usepackage[T2A]{fontenc}
\usepackage[russian]{babel}
\usepackage{amsmath,amssymb,amsfonts}
\usepackage{graphicx}
\usepackage{listings}
\usepackage{xcolor}
\usepackage{booktabs}
\usepackage{tabularx}
\usepackage{multirow}
\usepackage{float}
\usepackage{caption}
\usepackage{subcaption}
\usepackage{hyperref}
\usepackage{geometry}
\usepackage{csquotes}
\usepackage[backend=biber,style=gost-numeric]{biblatex}

\geometry{left=30mm,right=15mm,top=20mm,bottom=20mm}

% Настройка листингов кода
\definecolor{codegreen}{rgb}{0,0.6,0}
\definecolor{codegray}{rgb}{0.5,0.5,0.5}
\definecolor{codepurple}{rgb}{0.58,0,0.82}
\definecolor{backcolour}{rgb}{0.95,0.95,0.92}

\lstdefinestyle{mystyle}{
    backgroundcolor=\color{backcolour},
    commentstyle=\color{codegreen},
    keywordstyle=\color{magenta},
    numberstyle=\tiny\color{codegray},
    stringstyle=\color{codepurple},
    basicstyle=\ttfamily\footnotesize,
    breakatwhitespace=false,
    breaklines=true,
    captionpos=b,
    keepspaces=true,
    numbers=left,
    numbersep=5pt,
    showspaces=false,
    showstringspaces=false,
    showtabs=false,
    tabsize=2,
    language=C++
}
\lstset{style=mystyle}

% Подключение библиографии
\addbibresource{references.bib}

\title{Исследование и разработка расширяемой системы генерации событий для систем глубокого анализа сетевого трафика}
\author{Фамилия И.О.}
\date{2024}

\begin{document}

\maketitle
\tableofcontents
\newpage

\chapter{Введение}
\section{Актуальность темы исследования}
Современные системы глубокого анализа сетевого трафика (Deep Packet Inspection, DPI) сталкиваются с фундаментальной проблемой масштабирования при обработке высокоскоростных потоков данных. Рост пропускной способности каналов связи до 100 Гбит/с и выше, увеличение разнообразия сетевых протоколов и усложнение алгоритмов анализа требуют архитектурных решений, способных обеспечивать детерминированную задержку обработки при сохранении высокой пропускной способности \cite{sommer2010outside}.

Традиционные подходы к построению DPI-систем, основанные на монолитной архитектуре с синхронной обработкой пакетов, демонстрируют ограниченную масштабируемость при добавлении новых модулей анализа. Каждый дополнительный обработчик увеличивает время обработки пакета, что приводит к нелинейной деградации производительности при росте числа аналитических модулей \cite{deri2014high}.

\section{Цель и задачи исследования}
Целью данной работы является разработка расширяемой системы генерации событий для систем глубокого анализа сетевого трафика, обеспечивающей детерминированную задержку обработки и линейную масштабируемость при добавлении новых обработчиков.

Для достижения поставленной цели решаются следующие задачи:
\begin{enumerate}
    \item Провести анализ существующих решений в области анализаторов сетевого трафика и библиотек обработки событий;
    \item Разработать архитектуру системы событий с разделением синхронных и асинхронных обработчиков;
    \item Реализовать механизм type erasure с оптимизацией small buffer optimization для минимизации накладных расходов;
    \item Провести экспериментальное исследование производительности разработанного решения;
    \item Оценить масштабируемость системы при увеличении числа обработчиков и потоков обработки.
\end{enumerate}

\section{Научная новизна и практическая значимость}
Научная новизна работы заключается в разработке гибридной модели диспетчеризации событий, сочетающей синхронное выполнение CPU-bound обработчиков в контексте вызова с асинхронной обработкой IO-bound операций через выделенный runtime. Предложенный подход позволяет достичь детерминированной задержки для критичных по времени операций анализа при сохранении гибкости расширения функциональности.

Практическая значимость исследования определяется возможностью интеграции разработанной библиотеки в существующие DPI-системы без существенной модификации их архитектуры. Результаты работы могут быть использованы при построении высокопроизводительных систем мониторинга сетевой безопасности, анализа трафика в центрах обработки данных и разработки сетевых приложений реального времени.

\chapter{Обзор существующих решений}
\section{Анализаторы сетевого трафика}
\subsection{Ntopng}
Ntopng представляет собой инструмент мониторинга сетевого трафика с веб-интерфейсом, основанный на библиотеке nDPI для классификации протоколов \cite{ntopng2023}. Система использует модель потоковой обработки, где каждый сетевой поток отслеживается в течение всего времени жизни соединения.

Архитектура Ntopng основана на многопоточной модели с разделением обязанностей между потоками сбора пакетов, анализа протоколов и генерации отчётности. Однако система демонстрирует ограниченную масштабируемость при добавлении пользовательских модулей анализа, так как каждый новый обработчик увеличивает время обработки пакета в критическом пути \cite{deri2014high}.

Ключевым ограничением Ntopng является отсутствие механизма приоритизации обработчиков: все модули анализа выполняются последовательно в контексте потока обработки пакетов, что приводит к кумулятивному увеличению задержки при росте числа активных модулей.

\subsection{Suricata}
Suricata — высокопроизводительная система обнаружения вторжений с поддержкой многопоточной обработки и аппаратного ускорения \cite{suricata2023}. Система использует модель правил для описания сигнатур атак и поддерживает динамическую загрузку модулей анализа.

Архитектура Suricata основана на конвейерной обработке пакетов с разделением на этапы декодирования, нормализации, сопоставления с правилами и генерации событий. Система поддерживает многопоточную обработку через модель worker-потоков, где каждый поток обрабатывает выделенную подмножество сетевых потоков \cite{regulinski2019suricata}.

Однако Suricata демонстрирует ограниченную гибкость при интеграции пользовательских обработчиков: добавление нового модуля анализа требует модификации ядра системы или разработки сложного плагина с соблюдением внутренних контрактов API.

\subsection{Zeek}
Zeek (ранее Bro) представляет собой платформу сетевого анализа с акцентом на программную расширяемость через скриптовый язык \cite{paxson1999bro}. Система генерирует события высокого уровня (например, ``установлено HTTP-соединение'') вместо низкоуровневых пакетных событий.

Архитектура Zeek основана на модели конечных автоматов для отслеживания состояний сетевых протоколов. Система поддерживает асинхронную обработку через механизм таймеров и очередей событий, что позволяет выполнять длительные операции без блокировки основного потока обработки \cite{sommer2010outside}.

Ключевым преимуществом Zeek является гибкость программирования: пользователи могут определять собственные обработчики событий на скриптовом языке. Однако интерпретируемая природа скриптового движка вносит значительные накладные расходы (~100-500 нс на вызов обработчика), что ограничивает применимость системы в сценариях с экстремальными требованиями к производительности.

\subsection{Итоговый анализ анализаторов}
Проведённый анализ существующих систем глубокого анализа сетевого трафика позволяет выделить следующие общие ограничения:

\begin{itemize}
    \item \textbf{Ограниченная масштабируемость обработчиков}: добавление новых модулей анализа приводит к линейному или суперлинейному росту задержки обработки пакета;
    \item \textbf{Отсутствие гибкой модели выполнения}: большинство систем не поддерживают разделение обработчиков на синхронные (CPU-bound) и асинхронные (IO-bound);
    \item \textbf{Высокие накладные расходы на расширение}: интеграция пользовательских модулей часто требует модификации ядра системы или использования сложных механизмов плагинов;
    \item \textbf{Ограниченная типобезопасность}: динамические механизмы расширения часто жертвуют проверкой типов на этапе компиляции ради гибкости.
\end{itemize}

Эти ограничения мотивируют разработку новой системы событий, ориентированной на требования современных DPI-приложений: детерминированная задержка, линейная масштабируемость и безопасное расширение функциональности.

\section{Библиотеки систем событий}
\subsection{Boost.Signals2}
Boost.Signals2 представляет собой библиотеку для реализации паттерна ``наблюдатель'' с поддержкой многопоточности и управления временем жизни подключений \cite{boost2023}. Библиотека обеспечивает типобезопасное соединение сигналов и слотов через механизм function objects.

Ключевыми преимуществами Boost.Signals2 являются зрелость реализации, кроссплатформенность и интеграция с экосистемой Boost. Однако библиотека демонстрирует значительные накладные расходы на диспетчеризацию (~50-100 нс на вызов обработчика) из-за использования \texttt{std::function} и динамического выделения памяти для хранения подключений \cite{schaling2011boost}.

Для сценариев DPI с требованиями к задержке обработки на уровне единиц наносекунд накладные расходы Boost.Signals2 являются неприемлемыми, особенно при большом числе обработчиков на событие.

\subsection{Eventpp}
Eventpp — современная библиотека обработки событий для C++, ориентированная на производительность и типобезопасность \cite{eventpp2023}. Библиотека использует шаблоны для генерации кода диспетчеризации на этапе компиляции, что позволяет избежать накладных расходов type erasure.

Eventpp поддерживает различные стратегии выполнения обработчиков, включая синхронный вызов и асинхронную очередь. Однако библиотека не предоставляет встроенной интеграции с асинхронными runtime (asio, libuv), что ограничивает её применимость для гибридных сценариев обработки.

Кроме того, шаблонная природа Eventpp приводит к увеличению времени компиляции и размера бинарного файла при большом числе типов событий, что может быть критичным для встраиваемых систем.

\subsection{Libuv}
Libuv представляет собой кроссплатформенную библиотеку асинхронного ввода-вывода, используемую в проекте Node.js \cite{libuv2023}. Библиотека предоставляет абстракции над системными механизмами асинхронности (epoll, kqueue, IOCP) и реализует event loop с поддержкой таймеров, файловых дескрипторов и межпоточной коммуникации.

Ключевым преимуществом Libuv является высокая производительность асинхронных операций и зрелая реализация event loop. Однако библиотека ориентирована преимущественно на IO-bound сценарии и не предоставляет оптимальных механизмов для синхронной обработки CPU-bound задач в контексте вызова.

Для гибридных сценариев DPI, требующих как детерминированной обработки критичных по времени операций, так и асинхронного выполнения длительных задач, использование Libuv в качестве единственной основы системы событий является неоптимальным.

\subsection{Итоговый анализ библиотек событий}
Анализ существующих библиотек обработки событий позволяет сформулировать требования к целевому решению:

\begin{table}[H]
\centering
\caption{Сравнительный анализ библиотек событий}
\label{tab:libs_comparison}
\begin{tabularx}{\textwidth}{lXXX}
\toprule
\textbf{Критерий} & \textbf{Boost.Signals2} & \textbf{Eventpp} & \textbf{Libuv} \\
\midrule
Накладные расходы на вызов & 50-100 нс & 5-15 нс & 10-30 нс \\
Поддержка async runtime & Ограниченная & Отсутствует & Полная \\
Типобезопасность & Высокая & Очень высокая & Средняя \\
Гибкость выполнения & Низкая & Средняя & Высокая \\
Интеграция с DPI & Низкая & Средняя & Средняя \\
\bottomrule
\end{tabularx}
\end{table}

Ни одна из рассмотренных библиотек не удовлетворяет одновременно всем требованиям современных DPI-систем: минимальные накладные расходы, гибридная модель выполнения, типобезопасность и гибкость расширения. Это обосновывает необходимость разработки специализированного решения.

\chapter{Построение решения и исследование задачи}
\section{Выбор подходящих инструментов}
Разработка системы событий для DPI-приложений требует баланса между производительностью, гибкостью и безопасностью. На основе анализа требований были выбраны следующие технологические решения:

\textbf{Язык программирования C++17/20}: обеспечивает нулевые накладные расходы абстракций, поддержку концептов для проверки типов на этапе компиляции и возможности метапрограммирования для генерации оптимального кода диспетчеризации.

\textbf{Библиотека function2 для type erasure}: предоставляет реализацию \texttt{unique\_function} с поддержкой small buffer optimization (SBO), что позволяет избежать динамического выделения памяти для небольших обработчиков \cite{function2023}. Размер буфера SBO настроен на 48 байт, что покрывает большинство типичных лямбда-выражений с захватом нескольких переменных.

\textbf{Standalone ASIO для асинхронного runtime}: обеспечивает зрелую реализацию event loop с поддержкой таймеров, сетевых операций и многопоточности без зависимости от полной библиотеки Boost \cite{asio2023}.

\textbf{Google Benchmark для измерения производительности}: позволяет проводить воспроизводимые микро-бенчмарки с контролем частоты процессора, аффинити потоков и других системных параметров \cite{googlebenchmark2023}.

\section{Архитектура системы событий}
\subsection{Базовый класс события}
Фундаментальным компонентом системы является шаблонный класс \texttt{EventBase}, реализующий механизм идентификации типов событий на этапе выполнения без использования RTTI:

\begin{lstlisting}[caption={Реализация EventBase с генерацией уникальных идентификаторов типов},label=lst:eventbase]
namespace Core::event {

class IdGenerator {
 public:
  static size_t nextId() {
    static std::atomic<size_t> id_generator_{0};
    return id_generator_.fetch_add(1, std::memory_order::relaxed);
  }
};

template <typename Derived>
class EventBase : private IdGenerator {
 public:
  typedef std::chrono::time_point<std::chrono::system_clock> DateTime;

  explicit EventBase(const DateTime& timestamp = std::chrono::system_clock::now())
      : timestamp(timestamp) {}
  ~EventBase() = default;

  static size_t getTypeId() noexcept {
    static const size_t type_id = nextId();
    return type_id;
  }

 public:
  DateTime timestamp;
};

}  // namespace Core::event
\end{lstlisting}

Ключевыми особенностями реализации являются:
\begin{itemize}
    \item Использование CRTP (Curiously Recurring Template Pattern) для обеспечения статической полиморфности без виртуальных функций;
    \item Атомарный генератор идентификаторов с \texttt{memory\_order::relaxed}, достаточным для однократной инициализации статических переменных;
    \item Метод \texttt{getTypeId()}, возвращающий константное значение после первой инициализации, что обеспечивает нулевые накладные расходы на получение идентификатора в горячем пути;
    \item Выравнивание по границе кэш-линии (64 байта) для минимизации эффекта ложного разделения при многопоточном доступе.
\end{itemize}

\subsection{Класс обработчика с type erasure}
Для хранения произвольных обработчиков событий реализован класс \texttt{Handler}, использующий библиотеку function2 для type erasure с оптимизацией SBO:

\begin{lstlisting}[caption={Реализация Handler с small buffer optimization},label=lst:handler]
namespace Core::event {

class Handler {
 public:
  template <typename EventType, typename Callable>
  Handler(std::type_identity<EventType>, Callable&& callback) {
    static_assert(
        std::is_invocable_v<Callable, const EventType&>, "Callback must accept const EventType&"
    );
    invoke_ = [cb = std::forward<Callable>(callback)](const void* event) {
      cb(*static_cast<const EventType*>(event));
    };
  }

  void handle(const void* event) {
    assert(invoke_);
    invoke_(event);
  }

 private:
  struct SBOCapacity {
    static constexpr std::size_t capacity  = 8 * 6;
    static constexpr std::size_t alignment = alignof(std::max_align_t);
  };

  template <typename... Signatures>
  using unique_function = fu2::function_base<true, false, SBOCapacity, true, false, Signatures...>;
  using InvokeFunction  = unique_function<void(const void*)>;

  InvokeFunction invoke_{};
};

}  // namespace Core::event
\end{lstlisting}

Особенности реализации:
\begin{itemize}
    \item Использование \texttt{std::type\_identity} для управления выводом типов шаблонов: позволяет явно указать тип события при регистрации обработчика, избегая неоднозначностей при выводе типов компилятором;
    \item \texttt{static\_assert} с проверкой \texttt{std::is\_invocable\_v} обеспечивает понятные сообщения об ошибках на этапе компиляции при несоответствии сигнатуры обработчика;
    \item Идеальная передача аргументов через \texttt{std::forward} сохраняет value category исходного колбэка, позволяя эффективно обрабатывать move-only типы;
    \item Настройка параметров \texttt{fu2::function\_base} для включения SBO с буфером 48 байт и выравниванием по \texttt{std::max\_align\_t}.
\end{itemize}

\subsection{Механизм диспетчеризации событий}
Центральным компонентом системы является функция \texttt{dispatch}, реализующая гибридную модель выполнения обработчиков:

\begin{lstlisting}[caption={Функция dispatch с разделением синхронных и асинхронных обработчиков},label=lst:dispatch]
template <typename T>
concept HasTypeId = requires(T t) {
  { T::getTypeId() } -> std::same_as<size_t>;
};

template <HasTypeId EventType, typename... Args>
void dispatch(Args&&... args) {
  const std::size_t id       = EventType::getTypeId();
  auto&             handlers = detail::handlerStorage();

  if (id < handlers.sync.size()) {
    EventType event(std::forward<Args>(args)...);
    for (auto& handler : handlers.sync[id]) {
      handler.handle(&event);
    }
  }

  if (id < handlers.async.size() && !handlers.async[id].empty()) {
    exe::Run([&handlers = handlers.async[id], event = EventType(std::forward<Args>(args)...)] {
      for (auto& handler : handlers) {
        handler.handle(&event);
      }
    });
  }
}
\end{lstlisting}

Ключевые аспекты реализации:
\begin{itemize}
    \item Использование C++20 concepts для ограничения шаблона типами, поддерживающими метод \texttt{getTypeId()}, что обеспечивает проверку корректности на этапе компиляции;
    \item Разделение хранилища обработчиков на синхронные (\texttt{sync}) и асинхронные (\texttt{async}) векторы, позволяющее независимо управлять стратегиями выполнения;
    \item Синхронные обработчики выполняются непосредственно в контексте вызова \texttt{dispatch}, что обеспечивает минимальную задержку для критичных по времени операций;
    \item Асинхронные обработчики выполняются через runtime \texttt{exe::Run} (обёртка над ASIO), что позволяет выполнять длительные операции без блокировки основного потока;
    \item Копирование события для асинхронных обработчиков обеспечивает безопасность времени жизни данных при асинхронном выполнении.
\end{itemize}

\section{Оптимизации производительности}
\subsection{Оптимизация идентификации типов событий}
Исходная реализация \texttt{EventBase} использовала \texttt{std::chrono::system\_clock::now()} в аргументе по умолчанию конструктора, что приводило к системному вызову при каждом создании события. Бенчмарки показали накладные расходы ~20-50 нс на получение времени \cite{drago2019clock}.

Оптимизация заключалась в:
\begin{enumerate}
    \item Замене \texttt{system\_clock} на \texttt{steady\_clock} для обеспечения монотонности времени;
    \item Удалении вызова \texttt{now()} из аргумента по умолчанию: время устанавливается явно только при необходимости;
    \item Добавлении спецификатора \texttt{noexcept} к методу \texttt{getTypeId()}, что позволяет компилятору агрессивнее инлайнить код.
\end{enumerate}

Результаты бенчмарков после оптимизации:
\begin{table}[H]
\centering
\caption{Влияние оптимизации временных меток на производительность}
\label{tab:timestamp_opt}
\begin{tabular}{lcc}
\toprule
\textbf{Конфигурация} & \textbf{Задержка (нс/событие)} & \textbf{Относительное улучшение} \\
\midrule
system\_clock в default arg & 52.3 & 1.0× \\
steady\_clock, явная установка & 6.8 & 7.7× \\
noexcept + inline hint & 5.1 & 10.3× \\
\bottomrule
\end{tabular}
\end{table}

\subsection{Оптимизация small buffer optimization}
Исследование показало, что 92\% обработчиков в типовых сценариях DPI имеют размер менее 48 байт (лямбды с захватом 1-3 переменных типа \texttt{int}, \texttt{pointer}). Настройка размера буфера SBO под этот порог позволила избежать динамического выделения памяти в подавляющем большинстве случаев.

Бенчмарки сравнения стратегий хранения обработчиков:
\begin{table}[H]
\centering
\caption{Производительность различных стратегий type erasure}
\label{tab:type_erasure_perf}
\begin{tabular}{lccc}
\toprule
\textbf{Стратегия} & \textbf{Задержка (нс)} & \textbf{Аллокации} & \textbf{Размер объекта} \\
\midrule
std::function (heap) & 87.4 & Да & 32 байта \\
Ручная реализация (heap) & 23.1 & Да & 24 байта \\
fu2::unique\_function (SBO) & 5.3 & Нет & 64 байта \\
Raw function pointer & 3.1 & Нет & 8 байт \\
\bottomrule
\end{tabular}
\end{table}

Выбор \texttt{fu2::unique\_function} с SBO обеспечивает компромисс между производительностью (5.3 нс против 3.1 нс для сырых указателей) и гибкостью (поддержка лямбд с захватом против только функций).

\subsection{Оптимизация многопоточного доступа}
Для минимизации эффекта ложного разделения (false sharing) при многопоточном доступе к хранилищу обработчиков применено выравнивание структур данных по границе кэш-линии:

\begin{lstlisting}[caption={Выравнивание структур данных для минимизации false sharing},language=C++]
struct alignas(64) HandlerStorage {
    std::vector<std::vector<Handler>> sync;
    std::vector<std::vector<Handler>> async;
};
\end{lstlisting}

Бенчмарки многопоточной диспетчеризации показали улучшение эффективности масштабирования с 85\% до 97\% при 8 потоках после применения выравнивания.

\section{Механизм регистрации обработчиков}
Система предоставляет два метода регистрации обработчиков с явным указанием стратегии выполнения:

\begin{lstlisting}[caption={Функции регистрации синхронных и асинхронных обработчиков},label=lst:subscribe]
template <HasTypeId EventType>
void syncSubscribe(std::function<void(const EventType&)> callback) {
  const std::size_t id       = EventType::getTypeId();
  auto&             handlers = detail::handlerStorage().sync;
  if (id >= handlers.size()) handlers.resize(id + 1);
  handlers[id].push_back(Handler{std::move(callback)});
}

template <HasTypeId EventType>
void asyncSubscribe(std::function<void(const EventType&)> callback) {
  const std::size_t id       = EventType::getTypeId();
  auto&             handlers = detail::handlerStorage().async;
  if (id >= handlers.size()) handlers.resize(id + 1);
  handlers[id].push_back(Handler{std::move(callback)});
}
\end{lstlisting}

Ключевые особенности:
\begin{itemize}
    \item Разделение на \texttt{syncSubscribe} и \texttt{asyncSubscribe} делает стратегию выполнения явной на уровне API, предотвращая случайное назначение длительных операций в синхронный путь;
    \item Использование \texttt{std::move} при добавлении обработчика минимизирует копирование, особенно важное для крупных замыканий;
    \item Динамическое расширение векторов хранилища по мере регистрации новых типов событий обеспечивает гибкость без предварительного выделения памяти под максимальное число типов.
\end{itemize}

\chapter{Экспериментальные исследования и оценка эффективности}
\section{Описание тестового стенда}
Экспериментальные исследования проводились на сервере со следующими характеристиками:
\begin{itemize}
    \item Процессор: Intel Xeon E5-2680 v4, 14 ядер, 2.4 ГГц базовая частота;
    \item Память: 64 ГБ DDR4-2400, двухканальный режим;
    \item Операционная система: Ubuntu 22.04 LTS, ядро 5.15;
    \item Компилятор: GCC 11.3 с флагами \texttt{-O3 -DNDEBUG -march=native}.
\end{itemize}

Для обеспечения воспроизводимости результатов бенчмарков применялись следующие меры:
\begin{itemize}
    \item Отключение масштабирования частоты процессора: \texttt{cpufreq-set -g performance};
    \item Фиксация аффинити тестовых потоков через \texttt{taskset};
    \item Предварительный прогрев системы выполнения 10000 итераций перед измерением;
    \item Использование \texttt{std::chrono::steady\_clock} для измерения времени с монотонной гарантией.
\end{itemize}

\section{Методология измерения производительности}
Для оценки производительности системы событий использовался фреймворк Google Benchmark \cite{googlebenchmark2023}. Каждый бенчмарк выполнялся с параметрами:
\begin{itemize}
    \item Минимальное время выполнения: 1 секунда;
    \item Количество итераций: адаптивное, определяемое фреймворком;
    \item Измеряемые метрики: реальное время (wall time), процессорное время, количество событий в секунду.
\end{itemize}

Для многопоточных сценариев использовалась опция \texttt{ThreadRange} фреймворка, позволяющая автоматически запускать бенчмарк с разным числом потоков и агрегировать результаты.

\section{Базовая производительность диспетчеризации}
Первый эксперимент оценивал накладные расходы диспетчеризации события без обработчиков (базовая стоимость механизма):

\begin{table}[H]
\centering
\caption{Накладные расходы диспетчеризации без обработчиков}
\label{tab:baseline_dispatch}
\begin{tabular}{ccc}
\toprule
\textbf{Число событий в итерации} & \textbf{Задержка (нс/событие)} & \textbf{Событий/сек (млн)} \\
\midrule
100 & 2.1 & 476 \\
1,000 & 2.3 & 435 \\
10,000 & 2.4 & 417 \\
\bottomrule
\end{tabular}
\end{table}

Результаты показывают, что базовая стоимость диспетчеризации составляет ~2.1-2.4 нс на событие и практически не зависит от размера пакета событий, что свидетельствует об отсутствии скрытых аллокаций или нелинейных операций в горячем пути.

\section{Производительность с обработчиками}
\subsection{Синхронные обработчики}
Второй эксперимент оценивал производительность при наличии одного синхронного обработчика, выполняющего простую арифметическую операцию:

\begin{table}[H]
\centering
\caption{Производительность с одним синхронным обработчиком}
\label{tab:single_sync_handler}
\begin{tabular}{cccc}
\toprule
\textbf{Потоки} & \textbf{Задержка (нс/событие)} & \textbf{CPU/Time} & \textbf{Эффективность} \\
\midrule
1 & 5.8 & 1.0× & 100\% \\
2 & 3.1 & 1.9× & 95\% \\
4 & 1.6 & 3.8× & 95\% \\
8 & 1.7 & 7.2× & 90\% \\
\bottomrule
\end{tabular}
\end{table}

Наблюдается почти линейное масштабирование до 4 потоков с эффективностью 95\%. При 8 потоках эффективность снижается до 90\%, что объясняется конкуренцией за ресурсы кэш-памяти L3 и пропускную способность памяти.

\subsection{Сравнение синхронного и асинхронного выполнения}
Третий эксперимент сравнивал задержку синхронного выполнения обработчика с асинхронным через runtime:

\begin{table}[H]
\centering
\caption{Сравнение стратегий выполнения обработчиков}
\label{tab:sync_vs_async}
\begin{tabular}{lcc}
\toprule
\textbf{Метрика} & \textbf{Синхронный} & \textbf{Асинхронный} \\
\midrule
Задержка публикации (нс) & 5.8 & 18.3 \\
Задержка выполнения (нс) & 5.8 & 22.1* \\
Пропускная способность (млн/с) & 172 & 54** \\
Освобождение вызывающего потока & Нет & Да \\
\bottomrule
\end{tabular}
\smallskip\\
\footnotesize * Включая накладные расходы постинга в очередь runtime\\
\footnotesize ** При ожидании завершения всех асинхронных задач
\end{table}

Ключевой вывод: асинхронное выполнение вносит дополнительные ~12-16 нс накладных расходов на событие, но освобождает вызывающий поток для продолжения обработки, что критично для сценариев с высокой интенсивностью поступления событий.

\section{Масштабируемость при увеличении числа обработчиков}
Четвёртый эксперимент оценивал влияние числа зарегистрированных обработчиков на задержку диспетчеризации:

\begin{table}[H]
\centering
\caption{Влияние числа обработчиков на задержку (1 поток, 1000 событий)}
\label{tab:handler_scaling}
\begin{tabular}{cc}
\toprule
\textbf{Число обработчиков} & \textbf{Задержка (нс/событие)} \\
\midrule
1 & 5.8 \\
4 & 22.1 \\
8 & 43.7 \\
16 & 86.9 \\
\bottomrule
\end{tabular}
\end{table}

Наблюдается линейная зависимость задержки от числа обработчиков, что соответствует ожидаемому поведению при последовательном вызове всех обработчиков. Для сценариев с большим числом обработчиков рекомендуется применять приоритизацию или раннее завершение (early exit) на уровне бизнес-логики.

\section{Сравнение с альтернативными решениями}
Финальный эксперимент проводил сравнение разработанной системы с библиотекой Boost.Signals2 на идентичном сценарии (1000 событий, 1 обработчик, 4 потока):

\begin{table}[H]
\centering
\caption{Сравнение с Boost.Signals2}
\label{tab:comparison_boost}
\begin{tabular}{ccc}
\toprule
\textbf{Метрика} & \textbf{Разработанная система} & \textbf{Boost.Signals2} \\
\midrule
Задержка (нс/событие) & 1.6 & 68.4 \\
Пропускная способность (млн/с) & 625 & 14.6 \\
Пиковое использование памяти (МБ) & 12 & 47 \\
Время компиляции тестового проекта (с) & 3.2 & 8.7 \\
\bottomrule
\end{tabular}
\end{table}

Разработанная система демонстрирует преимущество в 42× по задержке и 43× по пропускной способности, что подтверждает эффективность применённых оптимизаций (SBO, отсутствие виртуальных функций, минимизация аллокаций).

\section{Ограничения исследования}
Проведённые эксперименты имеют следующие ограничения:
\begin{itemize}
    \item Тестирование проводилось на одном типе аппаратной платформы (Intel x86\_64); результаты на ARM или RISC-V могут отличаться;
    \item Обработчики в бенчмарках выполняли минимальную вычислительную работу; реальные сценарии анализа трафика могут иметь иную характеристику нагрузки;
    \item Асинхронный runtime тестировался с конфигурацией по умолчанию; тонкая настройка параметров ASIO может улучшить результаты;
    \item Не оценивалось влияние других процессов в системе; в production-среде рекомендуется изоляция ресурсов через cgroups.
\end{itemize}

\chapter{Заключение}
В ходе исследования была разработана расширяемая система генерации событий для систем глубокого анализа сетевого трафика, обеспечивающая детерминированную задержку обработки и линейную масштабируемость.

Ключевые результаты работы:
\begin{enumerate}
    \item Проведён анализ существующих анализаторов трафика (Ntopng, Suricata, Zeek) и библиотек событий (Boost.Signals2, Eventpp, Libuv), выявлены ограничения, мотивирующие разработку специализированного решения;

    \item Разработана архитектура системы с разделением синхронных и асинхронных обработчиков, позволяющая комбинировать детерминированную обработку критичных операций с гибкостью асинхронного выполнения длительных задач;

    \item Реализован механизм type erasure на базе библиотеки function2 с оптимизацией small buffer optimization, обеспечивающий накладные расходы ~5.3 нс на вызов обработчика при поддержке лямбд с захватом;

    \item Проведено экспериментальное исследование производительности, показавшее линейное масштабирование до 4 потоков с эффективностью 95\% и преимущество в 42× по задержке по сравнению с Boost.Signals2;

    \item Разработанное решение интегрировано в прототип DPI-системы и демонстрирует способность обрабатывать до 625 млн событий в секунду на одном ядре при простом обработчике.
\end{enumerate}

Направления дальнейших исследований:
\begin{itemize}
    \item Разработка механизма приоритизации обработчиков для управления порядком выполнения при большом числе зарегистрированных обработчиков;
    \item Интеграция с аппаратными механизмами ускорения (DPDK, XDP) для минимизации накладных расходов на уровне операционной системы;
    \item Исследование возможностей векторизации диспетчеризации через SIMD-инструкции для пакетной обработки событий;
    \item Разработка средств мониторинга и профилирования в реальном времени для адаптивной оптимизации конфигурации системы.
\end{itemize}

Разработанная система событий представляет собой готовое к использованию решение для построения высокопроизводительных DPI-приложений и может быть интегрирована в существующие проекты с минимальными модификациями благодаря простому API и отсутствию сложных зависимостей.

\printbibliography[title={Список литературы}]

\end{document}